\section{Physical units and the role of a continuum limit}
\label{sec:physical}

Choose a physical reference length $\ell_{\mathrm{ref}}>0$ and terminal
coupling $h=H$. A microscopic time separation of $n$ steps then equals
$a_Nn$, with $a_N=\ell_{\mathrm{ref}}L^{-N}$. Since
$e^{-\gap n}=e^{-(\gap/a_N)(a_Nn)}$, the mass in these units is
$m_{{\rm phys},N}=\gap(\beta_N(H))/a_N$. The proved bounds are
\begin{equation}
\label{eq:physical-bounds}
 \boxed{\displaystyle
  \frac{\log2}{M_0\ell_{\mathrm{ref}}}
  \le m_{{\rm phys},N}\le\frac{\log8}{\ell_{\mathrm{ref}}},
       \qquad N\ge K_0.}
\end{equation}
These are uniform bounds for cutoff theories. They require neither
convergence of their measures nor convergence of the numerical sequence
of masses. Equally, if one prescribes
$a(\beta)\Lambda=A\sqrt\beta e^{-\pi\beta}(1+o(1))$ with positive
fixed $A,\Lambda$, Theorem~\ref{thm:main} bounds the physical mass above
and below by positive multiples of $\Lambda$.

The half-logarithm in \eqref{eq:calibration} matters. Write
$\xi_{\rm lat}=\gap^{-1}$ for the inverse full gap. Knowing only
$\log\xi_{\rm lat}/\beta\to\pi$ would also allow
$\xi_{\rm lat}=\beta^{-1/2}e^{\pi\beta+\sqrt\beta}$, whose inverse
vanishes in these physical units. Our proof bounds the entire
multiplicative error.

\subsection{What would pass to an existing continuum theory}

We state separately the precise spectral implication if an appropriate
continuum limit is available. This is not used to prove
\eqref{eq:physical-bounds}.

\begin{proposition}
\label{prop:continuum}
Let $a_j\downarrow0$ and let $T_j$ be positive lattice transfer
contractions with normalized vacua $\Omega_j$ and full gaps $m_j$ satisfying
$\liminf_jm_j/a_j\ge m_*>0$. Suppose a limiting theory has Hilbert
space $\HH$, normalized vacuum $\Omega$, and a strongly continuous semigroup
$e^{-t\mathsf H}$ with $\mathsf H\ge0$ self-adjoint and
$\mathsf H\Omega=0$. Suppose a set of centered vectors $u$ dense in
$\Omega^\perp$ has
centered lattice approximants $u_j\perp\Omega_j$ and nonnegative
integers $n_j(t)$ such that
\[
 \norm{u_j}^2\to\norm u^2,\qquad
 \langle u_j,T_j^{n_j(t)}u_j\rangle
       \to\langle u,e^{-t\mathsf H}u\rangle,
 \qquad a_jn_j(t)\to t
\]
for every $t\ge0$. Then
\[
 \spec(\mathsf H)\subset\{0\}\cup[m_*,\infty),\qquad
 \ker\mathsf H=\C\Omega.
\]
\end{proposition}

\begin{proof}
For every $t>0$, the lattice gap gives
\[
 \langle u_j,T_j^{n_j(t)}u_j\rangle
 \le e^{-(m_j/a_j)a_jn_j(t)}\norm{u_j}^2.
\]
Pass to the limit and extend by density. For every $u\perp\Omega$,
$\langle u,e^{-t\mathsf H}u\rangle\le e^{-m_*t}\norm u^2$.
Positivity of its spectral measure excludes spectral weight below
$m_*$, including at zero, on that subspace.
\end{proof}

The norm convergence can instead be imposed after positive-time
regularization. More precisely, for each $u$ in a dense subset of
$\Omega^\perp$, suppose there are centered lattice vectors $v_j$ and
some $s>0$. Define their autocorrelations by
\[
 C_j(t)=\langle v_j,T_j^{\lfloor t/a_j\rfloor}v_j\rangle
 \qquad(t\ge0),
\]
and assume that
\[
 C_j(s+r)\longrightarrow\langle u,e^{-r\mathsf H}u\rangle
 \quad(r\ge0).
\]
The case $r=0$ gives $C_j(s)\to\|u\|^2$. For $t>s$, positivity of
the lattice spectral measure gives
$C_j(t)\le e^{-(m_j/a_j)(t-s+o(1))}C_j(s)$.
Apply this with $t=s+r$ and pass to the limit to obtain the same
quadratic-form bound as in the proposition. Thus no bound on
$\|v_j\|$, or on an unnormalized field renormalization constant,
is needed. Euclidean covariance is a separate condition if the energy
gap is to be interpreted as a relativistic mass gap.

An upper bound needs a surviving observable. Low-energy cutoff states
can disappear from a chosen limiting Hilbert space: the positive energy
measures $j^{-1}\delta_1+(1-j^{-1})\delta_R$, $R>1$, have spectral
bottom one but converge to $\delta_R$. A sufficient observable-survival condition is the following: in the
same limiting theory there are centered vectors $f,g$, with
$\|f\|,\|g\|\le1$, and a time separation $d\ge\ell_{\mathrm{ref}}$,
such that the separated cutoff block cross correlations converge to
$\langle f,e^{-d\mathsf H}g\rangle\ge1/8$. Then, writing
$m_{\mathrm{cont}}$ for the spectral bottom on $\Omega^\perp$,
\[
 1/8\le |\langle f,e^{-d\mathsf H}g\rangle|
 \le e^{-dm_{\mathrm{cont}}}
 \le e^{-\ell_{\mathrm{ref}}m_{\mathrm{cont}}},
\]
so $m_{\mathrm{cont}}\le(\log8)/\ell_{\mathrm{ref}}$.
The mixed-correlation convergence and norm bounds are additional
hypotheses; separate convergence of observables without these properties
would not suffice. These convergence requirements identify the
additional input for a statement about a continuum spectrum. The
finite and positive mass bounds in physical units have already been
proved without that input.
